\begin{enumerate}
    \item On dispose d'une solution $(S)$ de chlorure de potassium \ce{KCl} de
          concentration $c_1$ inconnue. Pour déterminer cette concentration
          plusieurs solutions de chlorure de potassium, de concentrations $c$,
          ont été réalisées. La mesure de la conductivité de ces solutions
          fournit les résultats suivants:
          % \begin{table}[H]
          %     \flushright
          %     \renewcommand{\arraystretch}{1.5}
          %     \begin{tabularx}{0.94\textwidth}{|p{3cm}|C|C|C|C|C|C|}
          %         \hline $c$~(\unit{\mmol\per\liter}) & $1$ & $2$ & $4$ &
          %                $6$ & $8$ & $10$ \\
          %         \hline $\sigma~(\times\qty{1e-4}{\siemens\per\meter})$ &
          %                149,8 & 299,6 & 599,2 & 898,7 & 1198,3 & 1498 \\
          %         \hline
          %     \end{tabularx}
          % \end{table}

          \begin{figure}[H]
              \centering
              \begin{tikzpicture}
                  \begin{axis}[
                        xmin=0,xmax=11.9,
                        ymin=0,ymax=1700,
                        width=9cm,height=7cm,
                        xlabel=$c~(\unit{\mmol\per\liter})$,
                        ylabel=$\sigma~(\times\qty{1e-4}{\siemens\per\meter})$,
                      ]
                      \addplot[
                          only marks,mark=+,mark options={very thick,scale=1.6},
                      ] table {data2.dat};
                      \addplot[domain=0:11,thick] {149.7937*x};
                  \end{axis}
              \end{tikzpicture}
          \end{figure}
          \begin{enumerate}
              \item Sachant que la conductivité de $(S)$ vaut~:
                    $\sigma_1=\qty{374,5e-4}{\siemens\per\meter}$, calculer
                    $c_1$ en \unit{\M}.
              \item Déduire des résultats expérimentaux la valeur de la
                    conductivité molaire $\lambda_{\ce{Cl-}}$ des ions chlorure.
          \end{enumerate}
    \item On ajoute à la solution $(S)$, sans variation de volume,
          $c_2=\qty{3e-3}{\M}$ de chlorure de sodium. Calculer la conductivité
          $\sigma^{\prime}$ de cette solution.
\end{enumerate}
% \begin{donnees}
%     Conductivités molaires des ions \ce{K$^+$_{(aq)}} et
%     \ce{Na$^+$_{(aq)}}~:
    
%     $\lambda_{\ce{K+}}=\qty{73.5e-4}{\lambda}$ et 
%     $\lambda_{\ce{Na+}}=\qty{50.1e-4}{\lambda}$
% \end{donnees}

